%% ma: L10-B13-0068
%% lop: 10
%% chuong: 5
%% bai: 13
%% dang: 10.13.13
%% loai: TL
%% muc: TH
%% dap_an: "-"
%% nguon: "chuyen-de-cac-so-dac-trung-cua-mau-so-lieu-khong-ghep-nhom-toan-10-chuong-trinh-moi.pdf (GV. Phan Nhật Linh) (P03), Bài 02 Dạng 2 Bài tập tự luận 5"
%% kien_thuc: "Lớp 10 Bài 13 (số trung bình, trung vị, tứ phân vị, mốt)"
%% trang_thai: chinh-thuc
%% ghi_chu: "Tự duyệt theo yêu cầu giáo viên 22/10/2026."
\begin{ex}
An lấy ra ngẫu nhiên $3$ quả bóng từ một hộp chứa nhiều bóng xanh và bóng đỏ, đếm xem có bao nhiêu bóng đỏ trong $3$ bóng đó rồi trả bóng lại hộp. An lặp lại phép thử trên $100$ lần và ghi lại kết quả như sau:
\begin{center}
\begin{tabular}{|c|c|c|c|c|}\hline
Số bóng đỏ & $0$ & $1$ & $2$ & $3$\\\hline
Số lần & $10$ & $30$ & $40$ & $20$\\\hline
\end{tabular}
\end{center}
Hãy tìm số trung bình, các tứ phân vị và mốt của mẫu số liệu.
\loigiai{$\overline x=\dfrac{0\cdot10+1\cdot30+2\cdot40+3\cdot20}{100}=1{,}7$.

Tần số tích lũy: $10;40;80;100$. $Q_2=\dfrac{x_{50}+x_{51}}2=2$; $Q_1=\dfrac{x_{25}+x_{26}}2=1$; $Q_3=\dfrac{x_{75}+x_{76}}2=2$.

Mốt $M_o=2$.}
\end{ex}
